\section{ОРГАНИЗАЦИЯ ВИРТУАЛЬНОЙ ВЫЧИСЛИТЕЛЬНОЙ СРЕДЫ}

\subsection{Выбор и развёртывание среды виртуализации}

В качестве среды виртуализации в данной работе использовался
программный продукт Oracle VM VirtualBox, выбор которого обусловлен
его свободным распространением, поддержкой широкого спектра
гостевых операционных систем и наличием официальной сборки для
архитектуры ARM64, что является необходимым условием для
развёртывания виртуальной среды на хост-системе с процессором Apple
Silicon (MacBook Air, чип Apple M2, 8 ГБ ОЗУ).

В соответствии с вариантом задания в качестве гостевой операционной
системы был выбран дистрибутив Fedora Server в редакции для
архитектуры aarch64. Установка и функционирование системы
осуществлялись в режиме интерфейса командной строки, без
графической оболочки.

\subsection{Архитектура виртуального кластера}

Для имитации вычислительных ресурсов суперкомпьютера было развёрнуто
пять виртуальных машин: один управляющий узел (\textit{control
node}), на котором размещаются диспетчер SLURM и средства разработки,
и четыре вычислительных узла (\textit{compute node}), предоставляющие
ресурсы для выполнения заданий. Конфигурация виртуальных машин,
выбранная с учётом ограничений хост-системы по объёму оперативной
памяти, приведена в таблице~\ref{tab:nodes_config}.

% TODO: проверьте и при необходимости скорректируйте конкретные цифры
\begin{table}[H]
\caption{Конфигурация узлов виртуального кластера}
\label{tab:nodes_config}
\centering
\begin{tabular}{|l|c|c|c|}
\hline
Узел & Роль & ЦП (ядра) & ОЗУ (МБ) \\
\hline
control & управляющий узел, slurmctld & 2 & 2048 \\
node1 & вычислительный узел, slurmd & 1 & 1024 \\
node2 & вычислительный узел, slurmd & 1 & 1024 \\
node3 & вычислительный узел, slurmd & 1 & 1024 \\
node4 & вычислительный узел, slurmd & 1 & 1024 \\
\hline
\end{tabular}
\end{table}

На каждом узле создана единая непривилегированная учётная запись
пользователя, используемая для запуска заданий и работы с диспетчером
SLURM; учётная запись \texttt{root} применялась исключительно для
установки и настройки системного программного обеспечения.

\subsection{Решение проблемы совместимости с архитектурой ARM}

В процессе установки гостевой операционной системы Fedora Server на
виртуальную машину с эмулируемым контроллером накопителей по
умолчанию (SATA/AHCI) была выявлена проблема несовместимости: на
хост-системе с архитектурой ARM (Apple M2) образ установки не
определял подключённый виртуальный диск, что приводило к невозможности
завершить установку.

Для устранения данной проблемы тип контроллера накопителя виртуальной
машины был изменён с SATA на \textbf{VirtIO SCSI} — паравиртуализированный
драйвер, обеспечивающий корректное взаимодействие гостевой системы с
диском при эмуляции архитектуры, отличной от архитектуры хост-машины.
После применения данного решения установка операционной системы
была успешно завершена на всех пяти виртуальных машинах.

% TODO: вставить скриншот настроек контроллера VirtualBox (Storage -> Controller: VirtIO SCSI)
% \begin{figure}[H]
% \centering
% \includegraphics[width=0.7\textwidth]{images/virtio_scsi.png}
% \caption{Настройка контроллера VirtIO SCSI в VirtualBox}
% \end{figure}

\subsection{Настройка виртуальной сети и разрешения имён}

Для обеспечения взаимодействия между узлами кластера был настроен
отдельный виртуальный сетевой адаптер в режиме
% TODO: уточнить режим сети — Host-only Adapter / Internal Network / NAT Network
\texttt{<режим сети VirtualBox>}, обеспечивающий изолированную
локальную сеть для всех пяти виртуальных машин. Каждому узлу был
присвоен статический IP-адрес в соответствии с таблицей~\ref{tab:ip}.

% TODO: подставить реальную сетевую маску/диапазон и адреса узлов
\begin{table}[H]
\caption{Адресация узлов виртуального кластера}
\label{tab:ip}
\centering
\begin{tabular}{|l|c|}
\hline
Имя узла & IP-адрес \\
\hline
control & 192.168.56.10 \\
node1 & 192.168.56.11 \\
node2 & 192.168.56.12 \\
node3 & 192.168.56.13 \\
node4 & 192.168.56.14 \\
\hline
\end{tabular}
\end{table}

Разрешение имён узлов обеспечено посредством редактирования файла
\texttt{/etc/hosts} на каждой виртуальной машине, что позволило
обращаться к узлам по символьным именам (\texttt{control},
\texttt{node1}--\texttt{node4}) без использования отдельного DNS-сервера.

\begin{lstlisting}[caption={Фрагмент файла /etc/hosts}]
192.168.56.10   control
192.168.56.11   node1
192.168.56.12   node2
192.168.56.13   node3
192.168.56.14   node4
\end{lstlisting}

\subsection{Настройка SSH с авторизацией по ключам}

На управляющем узле для непривилегированной учётной записи была
сгенерирована пара криптографических ключей (приватный и публичный)
с использованием утилиты \texttt{ssh-keygen}. Публичный ключ был
скопирован на все вычислительные узлы средствами утилиты
\texttt{ssh-copy-id}, что позволило выполнять удалённое подключение и
выполнение команд без запроса пароля.

\begin{lstlisting}[caption={Генерация и распространение ключей SSH}]
ssh-keygen -t ed25519 -C "cluster-key"
ssh-copy-id user@node1
ssh-copy-id user@node2
ssh-copy-id user@node3
ssh-copy-id user@node4
\end{lstlisting}

В конфигурационном файле демона SSH (\texttt{/etc/ssh/sshd\_config})
на всех узлах была отключена авторизация по паролю
(\texttt{PasswordAuthentication no}), что обеспечивает использование
исключительно ключевой авторизации.

\subsection{Установка и настройка munge}

Для аутентификации сообщений, которыми обмениваются демоны SLURM на
разных узлах, на управляющем узле был сгенерирован общий
секретный ключ \texttt{munge.key}, который был скопирован на все
вычислительные узлы с сохранением единых владельца и прав доступа
(пользователь \texttt{munge}). Служба \texttt{munge} была запущена и
добавлена в автозагрузку на всех пяти узлах кластера.

\begin{lstlisting}[caption={Развёртывание ключа munge}]
# На управляющем узле
sudo /usr/sbin/create-munge-key
scp /etc/munge/munge.key user@node1:/tmp/
# На каждом вычислительном узле
sudo mv /tmp/munge.key /etc/munge/munge.key
sudo chown munge: /etc/munge/munge.key
sudo chmod 400 /etc/munge/munge.key
sudo systemctl enable --now munge
\end{lstlisting}

\subsection{Установка и настройка SLURM}

На управляющем узле был установлен и настроен демон диспетчера
\texttt{slurmctld}, на вычислительных узлах — демон
\texttt{slurmd}. Основные параметры кластера определены в общем для
всех узлов конфигурационном файле \texttt{/etc/slurm/slurm.conf}:
имя кластера, адрес управляющего узла, список вычислительных узлов с
указанием количества процессоров и объёма оперативной памяти, а
также определение единственного раздела (\textit{partition}),
включающего все четыре вычислительных узла.

\begin{lstlisting}[caption={Фрагмент файла slurm.conf}]
ClusterName=virtual_cluster
SlurmctldHost=control
NodeName=node[1-4] CPUs=1 RealMemory=1024 State=UNKNOWN
PartitionName=debug Nodes=node[1-4] Default=YES MaxTime=INFINITE State=UP
\end{lstlisting}

После запуска демонов на всех узлах корректность конфигурации и
доступность вычислительных узлов проверялась командами
\texttt{sinfo} и \texttt{scontrol show nodes}.

% TODO: вставить скриншот вывода sinfo (все узлы в состоянии idle)
% \begin{figure}[H]
% \centering
% \includegraphics[width=0.9\textwidth]{images/sinfo_output.png}
% \caption{Вывод команды sinfo, подтверждающий готовность кластера}
% \end{figure}

\subsection{Клонирование вычислительных узлов}

После полной настройки и проверки работоспособности первого
вычислительного узла (установка ОС, сеть, SSH, munge, slurmd) данная
виртуальная машина использовалась как базовый образ для клонирования
оставшихся трёх вычислительных узлов средствами VirtualBox
(\textit{Clone} с режимом \textit{Full Clone}). После клонирования на
каждой копии выполнялись индивидуализация сетевых параметров (смена
IP-адреса и имени узла) и перегенерация машинного идентификатора
(\texttt{/etc/machine-id}) во избежание конфликтов в локальной сети.


